% from ega4-2-fr.tex
\oldpage[IV-2]{9}
est réduit (resp. intègre et intégralement clos). Si $f$ est fidèlement plat
et si $X$ est réduit (resp. normal), alors $Y$ est réduit (resp. normal).
\end{proposition}

Posons $\mathcal O_{f(x)}=A$, $\mathcal O_x=B$. Si $B$ est un $A$-module
plat, c'est aussi un $A$-module fidèlement plat $(0_{\mathrm I},6.6.2)$, donc
$A$ s'identifie à un sous-anneau de $B$~; si $B$ est réduit, il en est donc de
même de $A$. Supposons maintenant $B$ intègre et intégralement clos, et soit
$L$ son corps des fractions~; alors $A\subset B$ est intègre~; désignons par
$K\subset L$ son corps des fractions. L'hypothèse entraîne que
$B\cap K=A$ (Bourbaki, \emph{Alg. comm.}, chap.~I, \S~3,
n\textsuperscript{o}~5, prop.~10). Si alors $t\in K$ est entier sur $A$, il
est aussi entier sur $B$, donc appartient à $B$ par hypothèse, et par suite
$t\in A$, ce qui prouve que $A$ est intégralement clos.

\begin{proposition}[2.1.14]
\label{IV.2.1.14-fr}
Soit $f:X\to Y$ un morphisme fidèlement plat. Si $X$ est localement intègre
et si l'espace sous-jacent à $Y$ est localement noethérien, alors $Y$ est
localement intègre.
\end{proposition}

En effet, tout $y\in Y$ est de la forme $f(x)$ pour un $x\in X$, et par
hypothèse $\mathcal O_y$ s'identifie à un sous-anneau de $\mathcal O_x$
$(0_{\mathrm I},6.6.1)$~; comme $\mathcal O_x$ est intègre, il en est de même
de $\mathcal O_y$ et cela prouve la proposition (I, 5.1.4).

\subsection{Modules fidèlement plats sur les préschémas}
\label{subsection:IV.2.2-fr}

\begin{proposition}[2.2.1]
\label{IV.2.2.1-fr}
Soient $f:X\to Y$ un morphisme de préschémas, $\mathscr F$ un
$\mathcal O_X$-Module quasi-cohérent. Les conditions suivantes sont
équivalentes~:
\begin{enumerate}
\item[\emph{a)}] Pour tout changement de base $g:Y'\to Y$, si l'on pose
$X'=X\times_Y Y'$, le foncteur
$\mathscr G'\mapsto\mathscr F\otimes_Y\mathscr G'$, de la catégorie des
$\mathcal O_{Y'}$-Modules quasi-cohérents dans celle des
$\mathcal O_{X'}$-Modules quasi-cohérents, est exact et fidèle.
\item[\emph{a$'$)}] La condition \emph{a)} est vérifiée pour tous les
morphismes canoniques $g:\operatorname{Spec}(\mathcal O_y)\to Y$ (I, 2.4.1),
où $y$ parcourt $Y$.
\item[\emph{a$''$)}] La condition \emph{a)} est vérifiée pour toutes les
immersions canoniques $Y'\to Y$, où $Y'$ parcourt l'ensemble des
sous-schémas ouverts affines de $Y$.
\item[\emph{b)}] $\mathscr F$ est $Y$-plat et, pour tout $y\in Y$, si l'on
désigne par $X_y$ la fibre $f^{-1}(y)$, le $\mathcal O_{X_y}$-Module
$\mathscr F_y=\mathscr F\otimes_Y k(y)$ est $\neq0$.
\end{enumerate}
\end{proposition}

Il est clair que \emph{a)} implique \emph{a$'$)} et \emph{a$''$)}~; la
condition \emph{a$'$)} implique d'abord que $\mathscr F$ est $Y$-plat
(2.1.3)~; d'autre part elle implique que pour tout $y\in Y$, le foncteur
$N\mapsto\mathscr F\otimes_{\mathcal O_y}\widetilde N$ est fidèle dans la
catégorie des $\mathcal O_y$-modules~; prenant en particulier $N=k(y)$, on
obtient la seconde assertion de \emph{b)}. Pour montrer que \emph{b)}
implique \emph{a)}, on peut se limiter au cas où $Y$ est \emph{affine}, la
question étant locale sur $Y$. De même, pour prouver que \emph{a$''$)}
implique \emph{a)}, on est ramené à prouver que lorsque $Y$ est affine, le
fait que $\mathscr G\mapsto\mathscr F\otimes_Y\mathscr G$ soit un foncteur
exact et fidèle entraîne la condition \emph{a)}. Autrement dit, on est ramené
à prouver la proposition plus précise suivante~:

\begin{proposition}[2.2.2]
\label{IV.2.2.2-fr}
Soient $Y=\operatorname{Spec}(A)$ un schéma affine, $f:X\to Y$ un morphisme
de préschémas, $\mathscr F$ un $\mathcal O_X$-Module quasi-cohérent. La
condition \emph{a)} de (2.2.1) est équivalente à chacune des suivantes~:
\begin{enumerate}
\item[\emph{b$'$)}] $\mathscr F$ est $Y$-plat et, pour tout point fermé $y$
de $Y$, on a $\mathscr F_y\neq0$.
